\documentclass[unicode,10pt]{beamer}


\input{~/Documents/y-ta-san.github.io/lt_head/details/pkgimport_v2.0_beamer.tex}
%%%%%%%%%%%%%%%%%%%%%%%%%%%%%%%%%%%%%%%%%%%%%%%%%%
% redefinition of command
% history 
% v1.0 : 27 Feb, 2026
% v2.0 : 24 Jun, 2026
%%%%%%%%%%%%%%%%%%%%%%%%%%%%%%%%%%%%%%%%%%%%%%%%%%

% kana counter (using expl3)
\ExplSyntaxOn
\makeatletter
\cs_new:Npn \kana_map:n #1 {
  \ifcase #1 \or ア \or イ \or ウ \or エ \or オ \or カ \or キ \or ク \or ケ \or コ
  \or サ \or シ \or ス \or セ \or ソ \or タ \or チ \or ツ \or テ \or ト 
  \or ナ \or ニ \orヌ \or ネ \or ノ \or ハ \or ヒ \or フ \or ヘ \or ホ
  \or マ \or ミ \or ム \or メ \or モ \or ヤ \or ユ \or ヨ \or ラ \or リ
  \or ル \or レ \or ロ \or ワ \or ヲ \or ン \fi
}
\newcommand{\kana}[1]{
  \expandafter \kana_map:n \csname c@#1\endcsname
}
\makeatother
\ExplSyntaxOff

\renewcommand{\emph}[1]{{\sffamily\bfseries\gtfamily\kanjiseries{m}\selectfont #1}}

%%% mathematics
% change font
\newcommand{\mr}[1]{\symup{#1}}
\newcommand{\mb}[1]{\symbf{#1}}
\newcommand{\bs}[1]{\symbf{#1}}
\newcommand{\bb}[1]{\symbb{#1}}
\newcommand{\scr}[1]{\symscr{#1}}

% sets
\newcommand{\R}{\symbb{R}}
\newcommand{\N}{\symbb{N}}
\newcommand{\C}{\symbb{C}}
\newcommand{\Z}{\symbb{Z}}

% const.
\renewcommand{\i}{\symup{i}}
\newcommand{\e}{\symup{e}}

% matrix, average, order
\newcommand{\mqty}[1]{\begin{matrix} #1 \end{matrix}}
\newcommand{\avg}[1]{\langle #1 \rangle}
\newcommand{\order}[1]{\mathcal{O}(#1)}
\newcommand{\eval}[1]{
  #1 \mathclose{
    \rule[-2.0ex]{0.4pt}{5.5ex}\,
  }
}

% differential op
\newcommand{\grad}{\nabla}
\renewcommand{\div}{\nabla \cdot}
\newcommand{\rot}{\nabla \times}
\newcommand{\nb}{\nabla}

%% DIFFCOEFF PKG SETTING
\difdef { l } { dn } { style = d^ }
\newcommand{\dn}[2]{ \dl.dn.[#1]{#2}\, }
\difdef{l}{p}{op-symbol=\partial}
\newcommand{\dlp}{ \dl.p. }
\difdef { l } {} {outer-Rdelim = \,}
\difdef {f, s} {D}{op-symbol = \symup{D}}
\newcommand{\Dl}[1]{\upDelta{#1}}

%%% abbrv.
\newcommand{\eps}{\varepsilon}

%%% layout
\newcommand{\qq}{\quad}
\newcommand{\qqtext}[1]{\quad \text{#1} \quad}
\newcommand{\ds}{\displaystyle}

%%% color
\definecolor{mint}{HTML}{3EB370}
\definecolor{dark-mint}{HTML}{2E8555}
\renewcommand{\CancelColor}{\color{ForestGreen}}
\newcommand{\blue}[1]{\textcolor{blue}{#1}}
\newcommand{\red}[1]{\textcolor{red}{#1}}

%%% other configuration
\AtBeginDocument{
  \renewcommand{\figurename}{Figure\,}
  \renewcommand{\proofname}{\bfseries\sffamily\gtfamily Proof.}
}

% equation number (section.eq)
\numberwithin{equation}{section}



\usepackage{geometry}
\geometry{papersize={160mm,90mm}}
\usepackage{luatexja}
\usepackage[haranoaji]{luatexja-preset}
%\usepackage{unicode-math}
%\usepackage{diffcoeff}

\usetheme{default}
%\setmathfont{ArsenalMath}
%\setmathfont{kpmath-sans.otf}
\setmathfont{Concrete Math}
%\setmathfont{Concrete Math}
%\setmainfont[
%  BoldFont={Arsenal Bold},
%  BoldItalicFont={Arsenal BoldItalic},
%  ItalicFont={Arsenal Italic}
%]{Arsenal Regular}
%\setmainfont{kpsans-regular.otf}
\setmainfont{CMU Concrete}
\setmainjfont[
  BoldFont={A P-OTF UD ShinGoCon90 Pr6N M},
]{A P-OTF UD ShinGoCon90 Pr6N R}

% 配色の設定
\definecolor{LightOrange}{HTML}{FFE66F}
\definecolor{AccentOrange}{HTML}{FF8C00}

\setbeamercolor{background canvas}{bg=white}
\setbeamercolor{frametitle}{bg=LightOrange, fg=black!90}
\setbeamercolor{title}{bg=LightOrange, fg=black!90}
\setbeamercolor{itemize item}{fg=AccentOrange}
\setbeamercolor{itemize subitem}{fg=AccentOrange}

\setbeamercolor{footline}{bg=LightOrange, fg=black!90}

\setbeamertemplate{navigation symbols}{}

\makeatletter
\setbeamertemplate{frametitle}{
  \nointerlineskip
  \begin{beamercolorbox}[wd=\paperwidth, leftskip=0.4cm, rightskip=0.4cm]{frametitle}
    \vskip 0.1\zh
    {\scriptsize\color{black!60}%
     \strut\insertsection
     \ifx\insertsubsection\@empty\else\ > \insertsubsection\fi
     \par}
    \vspace{-0.2\zh}
    {\large\textbf{\insertframetitle}}
    \vskip 0.3\zh
  \end{beamercolorbox}
}
\makeatother
\setbeamertemplate{footline}{
  \nointerlineskip
  \begin{beamercolorbox}[wd=\paperwidth, ht=1.35\zh, dp=.75\zh, leftskip=4\zh, rightskip=4\zh]{footline}
    {\scriptsize \hfill \insertframenumber{} / \inserttotalframenumber}
  \end{beamercolorbox}

\definecolor{AlertRed}{HTML}{FF6B6B}
\definecolor{ExampleGreen}{HTML}{88D49E}

% 通常のblock
\setbeamertemplate{blocks}[default]
\setbeamercolor{block title}{bg=LightOrange, fg=black!90}
\setbeamercolor{block body}{bg=LightOrange!20, fg=black!90}

% 警告・強調用のalertblock
\setbeamercolor{block title alerted}{bg=AlertRed, fg=white}
\setbeamercolor{block body alerted}{bg=AlertRed!10, fg=black!90}

% 具体例用のexampleblock
\setbeamercolor{block title example}{bg=ExampleGreen, fg=black!90}
\setbeamercolor{block body example}{bg=ExampleGreen!20, fg=black!90}
}


%\usepackage[most]{tcolorbox}

% block用追加カラーの定義（前回と同じ）
\definecolor{AlertRed}{HTML}{FF6B6B}
\definecolor{ExampleGreen}{HTML}{88D49E}

% tcolorboxの共通スタイル設定
\tcbset{
  my_beamer_block_style/.style={
    enhanced,
    boxrule=0pt, % 枠線をなしにしてフラットに
    coltitle=black!90, % タイトルの文字色
    fonttitle=\bfseries, % タイトルを太字に
    sharp corners, % 四角いデザイン
    toptitle=1.5mm, bottomtitle=1.5mm, % タイトル部分の上下余白
    left=3mm, right=3mm, top=2mm, bottom=2mm, % 本文部分の余白
  }
}

% 独自のblock環境（myblock）を新しく定義
\newtcolorbox{box1}[1]{
  boxrule=0pt,
  my_beamer_block_style,
  colbacktitle=LightOrange,
  colback=LightOrange!20,
  width=0.8\linewidth,
  toptitle=0.5mm, bottomtitle=0.5mm,
  center,
  title=#1
}
